\chapter{Thalès et triangles semblables}
\label{exercices:14}
\section*{Exercice 14.1 -- Configuration emboîtée}
Dans $ABC$, $M\in[AB]$, $N\in[AC]$, $(MN)\parallel(BC)$. On donne $AB=15$, $AC=12$, $BC=18$, $AM=10$ en centimètres. Calculer $AN$, $MN$ et $MB$.

\section*{Exercice 14.2 -- Prouver un parallélisme}
Les points $M$ et $N$ sont respectivement sur $[AB]$ et $[AC]$, avec $AM=7$, $AB=10$, $AN=10{,}5$, $AC=15$. Prouver que $(MN)\parallel(BC)$. Que conclurait-on si $AN=11$ ?

\section*{Exercice 14.3 -- Configuration papillon}
Les points $B,A,M$ sont alignés dans cet ordre, de même que $C,A,N$. On sait $(MN)\parallel(BC)$, $AM=6$, $AB=9$, $AC=6$, $BC=12$. Calculer $AN$, $MN$, $BM$ et $CN$.

\section*{Exercice 14.4 -- Angles et similitude}
Deux triangles ont respectivement des angles $(40\degre,60\degre,80\degre)$ et $(60\degre,80\degre,40\degre)$. Le côté opposé à $40\degre$ mesure $6$ dans le premier et $9$ dans le second. Le premier a un périmètre de $24$ et une aire de $18$. Calculer le périmètre et l'aire du second.
